\documentclass[11pt]{article}
\usepackage[utf8]{inputenc}
\usepackage{amsmath,amssymb,amsthm}
\usepackage[margin=1in]{geometry}
\usepackage{hyperref}
\usepackage{xcolor}
\usepackage{enumitem}
\usepackage{newunicodechar}
\newunicodechar{ℝ}{\ensuremath{\mathbb{R}}}
\newunicodechar{ℕ}{\ensuremath{\mathbb{N}}}
\newunicodechar{⟨}{\ensuremath{\langle}}
\newunicodechar{⟩}{\ensuremath{\rangle}}
\newunicodechar{Γ}{\ensuremath{\Gamma}}
\newunicodechar{·}{\ensuremath{\cdot}}
\newunicodechar{≤}{\ensuremath{\le}}
\newunicodechar{→}{\ensuremath{\to}}

\newcommand{\gptbox}[1]{\par\medskip\begin{quote}\small\textbf{\textcolor{green!50!black}{GPT-5.5 Pro:}} #1\end{quote}\medskip}

\title{Tide retrospective:\\\emph{mixed even-power 2D moments
(quartic-sextic potential)}}
\author{Claude Opus 4.7 (1M ctx), with Daniel Murfet}
\date{2026-05-07}

\begin{document}
\maketitle

\begin{abstract}
For the 2D separable potential $L(x, y) = x^4/24 + y^6/720$ at
temperature $t > 0$ and any $j, k : \mathbb N$, the joint even-power
moment $\langle x^{2j} \cdot y^{2k} \rangle_{2D, t}$ has the closed
form
$$
  \left( (24/t)^{j/2} \cdot \tfrac{\Gamma((2j+1)/4)}{\Gamma(1/4)} \right)
  \cdot
  \left( (720/t)^{k/3} \cdot \tfrac{\Gamma((2k+1)/6)}{\Gamma(1/6)} \right).
$$
This is the obvious closed-form sequel to the just-merged A2 tide
(quartic-sextic mixed potential infrastructure). New file
\texttt{Laplace/TwoD/QuarticSexticMoments.lean}, 57~lines (proof body
$\sim 8$ lines; the rest is doc-comment); zero \texttt{sorry}, zero
\texttt{axiom}, zero \texttt{native\_decide}. Single-attempt clean
build via two \texttt{rw} calls. Fourth tide today (after I4, N1, Q1)
where the Step~2 advisory carried sufficient tactical detail to make
Step~3 mechanical.
\end{abstract}

\section{Setting}

A2 (\texttt{tide/quartic-sextic-mixed}, just landed) added the 2D
separable-potential infrastructure for the quartic-sextic case. Its
key public theorem reduces the 2D Gibbs expectation of a separable
observable $f(x) g(y)$ to a product of the two 1D Gibbs expectations
for the quartic and sextic potentials, given integrability of each
piece against its 1D Boltzmann weight.

The closed-form 1D even moments
$\langle x^{2n} \rangle_{\text{quartic}, t}$ and
$\langle x^{2n} \rangle_{\text{sextic}, t}$ have been formalised since
the early Tide steps. So the joint even-moment closed form
$\langle x^{2j} y^{2k} \rangle_{2D, t}$ is essentially a one-line
algebraic combination of three already-formalised pieces. A2's
retrospective explicitly named this as the natural sequel.

The tide-pick v5 survey ranked it as the top recommendation. The
Step~2 deliberation between Claude and GPT-5.5~Pro converged in a
single round on Candidate~A (just the closed-form combined moment);
GPT additionally flagged that the existing public factor theorem
already exposes the product-structure form, so a separate
\texttt{pow\_pow\_factor} intermediate would be redundant.

The markdown tide log\footnote{In the SRI repo at
\texttt{projects/primer/tide-log/}, file
\texttt{2026-05-07-tide-mixed-even-2d-moments.md}.} carries the
deliberation; the GPT consult is preserved at the sibling
\texttt{gpt\_responses\_t1/t1\_v1.md}.

\section{The thing formalised}

\begin{quote}\itshape
For the 2D separable potential $L(x, y) = x^4/24 + y^6/720$, any
$t > 0$, and any non-negative integers $j, k$,
$$
  \langle x^{2j} \cdot y^{2k} \rangle_{2D, t}
  = \left( (24/t)^{j/2} \cdot \frac{\Gamma((2j+1)/4)}{\Gamma(1/4)} \right)
    \cdot \left( (720/t)^{k/3} \cdot \frac{\Gamma((2k+1)/6)}{\Gamma(1/6)} \right).
$$
\end{quote}

\noindent The Lean signature\footnote{\texttt{Laplace/TwoD/%
QuarticSexticMoments.lean}, theorem
\texttt{gibbsExpectation\_quarticSextic\_pow\_pow\_eq}.}:

\begin{verbatim}
theorem gibbsExpectation_quarticSextic_pow_pow_eq
    (j k : ℕ) {t : ℝ} (ht : 0 < t) :
    gibbsExpectation
        (addSeparable Laplace.OneD.quarticPotential
                      Laplace.OneD.sexticPotential) t
        (fun z : ℝ × ℝ => z.1 ^ (2 * j) * z.2 ^ (2 * k))
      = ((24 / t) ^ ((j : ℝ) / 2)
            * Real.Gamma ((2 * j + 1 : ℝ) / 4)
            / Real.Gamma ((1 : ℝ) / 4))
        * ((720 / t) ^ ((k : ℝ) / 3)
            * Real.Gamma ((2 * k + 1 : ℝ) / 6)
            / Real.Gamma ((1 : ℝ) / 6))
\end{verbatim}

\noindent The 2D space is $\mathbb{R} \times \mathbb{R}$ (per A2's
convention), not $\mathrm{Fin}\,2 \to \mathbb R$; observables index
into it via \texttt{z.1} and \texttt{z.2}. GPT specifically called
this out at deliberation as a likely point of confusion otherwise.

The closed-form RHS is kept as a literal product of the two 1D
forms — no attempt to combine $(24/t)^{j/2}$ and $(720/t)^{k/3}$
into a single $t$-power. GPT flagged that combining "creates
gratuitous \texttt{rpow}/positivity/algebra overhead" with no
mathematical or readability gain.

\section{Proof strategy}

The proof is two \texttt{rw} calls:

\begin{enumerate}[leftmargin=1em,itemsep=0.2em]
\item Apply A2's separable-observable factor theorem with
  $f := x \mapsto x^{2j}$ and $g := y \mapsto y^{2k}$, using the
  existing 1D integrability-of-monomial-against-Boltzmann-weight
  witnesses for the quartic and sextic potentials. The 2D
  expectation becomes a product of two 1D expectations.
\item Substitute the 1D even-moment closed forms (one quartic, one
  sextic) to obtain the Gamma-form RHS.
\end{enumerate}

\noindent The proof body is eight lines of Lean.

\section{Roadblocks and resolutions}

\emph{None}. Single-attempt clean build, zero warnings. This is the
fourth tide today (after I4, N1, Q1) where a clean Step~2 advisory
made Step~3 mechanical. The pattern is now reliable: when the Step~2
deliberation correctly identifies the integrability bundle, the
namespace conventions, and the verbatim form of the closed-form
RHS, the proof writes itself.

\gptbox{Big Lean caveat: copy the RHS verbatim from the 1D theorems.
In particular, match exactly whether they use \texttt{\^{}} or
\texttt{Real.rpow}, \texttt{((2 * j + 1 : ℝ) / 4)} vs extra casts,
parenthesisation of \texttt{*} and \texttt{/}. If you match the 1D
theorem statements exactly, the proof is basically two rewrites. If
you "prettify" the RHS, you'll buy yourself avoidable algebra.}

\noindent The implementation took this guidance literally — the RHS
is character-for-character a product of the two 1D theorem RHSes.
The proof is two rewrites.

\paragraph{Numerical sanity pre-Step~3.} Eight cases at
$t \in \{1, 2\}$, $(j, k) \in \{0, 1, 2\}^2$, all agree to
${\sim}10^{-13}$ relative error against \texttt{scipy.integrate.dblquad}.
The identity is exact (a structural separability fact for the 2D
measure), not asymptotic.\footnote{Script preserved at
\texttt{projects/primer/tide-log/gpt\_responses\_t1/t1\_numerics.py}.}

\section{What was Mathlib, what was new}

\paragraph{Mathlib provided.} Nothing directly — this tide is
entirely above the Mathlib stack, composing already-formalised
Laplace/laplace pieces.

\paragraph{What was new.} A single public theorem
(\texttt{gibbsExpectation\_quarticSextic\_pow\_pow\_eq}) that
combines three already-formalised pieces:

\begin{itemize}[leftmargin=1em,itemsep=0.2em]
\item A2's separable-observable factor theorem (the 2D-to-1D
  reduction);
\item the existing 1D integrability witnesses for monomials against
  the quartic and sextic Boltzmann weights;
\item the 1D even-moment closed forms for the quartic and sextic
  potentials.
\end{itemize}

\noindent The novelty is just the composition; no new lemma machinery.

\section{Lessons}

\begin{itemize}[leftmargin=1em,itemsep=0.2em]
\item \emph{Sometimes the right Tide is the trivial composition.}
  A2 just landed; the obvious sequel is a 30-line composition
  theorem. There is no value-add in proving anything more elaborate
  in this Tide step. Resist the impulse to bundle further
  generalisations (odd-power vanishing, generic pure-monomial pairs,
  etc.) — those are separate Tide steps.

\item \emph{The "match RHS verbatim" rule.} Whenever a Tide composes
  two existing theorems with literal product-of-RHS targets, write
  the headline RHS character-for-character as the product of the
  two 1D RHSes. This makes the proof a pure \texttt{rw} chain. Any
  algebraic re-expression — even a \emph{semantically equivalent}
  one — will introduce avoidable \texttt{rpow} and positivity steps.

\item \emph{The "is the existing public theorem already enough?"
  test.} GPT correctly noted that proposing a separate
  \texttt{pow\_pow\_factor} public intermediate (Candidate~B) would
  be redundant given A2's existing \texttt{*\_factor} theorem
  already exposes the product structure for arbitrary $f, g$. When a
  Tide proposes a sub-result that's already a special case of an
  existing public theorem, the right answer is "the existing theorem
  is enough; just show how to use it". Worth recording for future
  Tide deliberations.
\end{itemize}

\section{Follow-ups}

\begin{itemize}[leftmargin=1em,itemsep=0.2em]
\item \emph{Mixed odd-power vanishing} (Candidate~C, deferred):
  $\langle x^{2j+1} y^{2k} \rangle_{2D, t} = 0$ and the symmetric
  variant. Each follows from the factor theorem plus the existing
  1D odd-moment vanishing theorems. About ${\sim}30$ lines for the
  parity-completion pair.

\item \emph{Other 2D mixed pure-monomial pairs.} Quartic-quartic,
  sextic-sextic, $x^{2j}/((2j)!)$ paired with $y^{2k}/((2k)!)$ for
  arbitrary $(j, k)$. Each is a 30-line corollary of A3 once the
  named 1D pure-monomial families exist; A2's retrospective named
  this as a follow-up. Best deferred until A1 (the generic
  even-monomial 1D template) lands.

\item \emph{Generic-$k$ J-function-style 2D moment asymptotic.} The
  natural \emph{asymptotic} sequel: as $t \to \infty$, the 2D
  moment behaves as the product of the leading-order 1D scalings.
  Quartic side scales as $t^{-(2j+1)/4}$, sextic side as
  $t^{-(2k+1)/6}$, joint scaling $t^{-(2j+1)/4 - (2k+1)/6}$.
  Tide-shaped; ${\sim}80$ lines combining
  \texttt{kth\_jfunction\_asymptotic} machinery from N1.

\item \emph{\texttt{\textbackslash leanref}-tag the primer's 2D
  examples} once the primer has TeX statements for any of the 2D
  moment formulas. Pin to this tide's commit. Process tide; pairs
  with K2.
\end{itemize}

\section*{Acknowledgements}

GPT-5.5~Pro contributed the deliberation-stage advisory at Step~2:
the verdict for Candidate~A over B (B is redundant given A2's
existing public factor theorem) and C (better as a separate
parity-completion follow-up); the \texttt{z.1}/\texttt{z.2} naming
guidance; and the verbatim-RHS rule that made the proof a pure
two-line \texttt{rw} chain. The full consult is preserved alongside
the markdown tide log.\footnote{In the SRI repo at
\texttt{projects/primer/tide-log/gpt\_responses\_t1/}, file
\texttt{t1\_v1.md}.} The proof landed single-attempt with no
mid-proof thrashing-rescue consult needed.

\end{document}
